%%%%%%%%%%%%%%%%%%%%%%%%%%%%%%%%%%%%%%%%%%%%%%%%%%%%%%%%%%%%%%%%%%%%%%%%%%%%%%%%
% Revisión Sistemática de la Literatura (RSL) — Clasificación de requerimientos
% no funcionales con desbalance de clases: aumentación de datos, BERT y LLM.
% Formato: IEEEtran (conference). Bibliografía: referencias_nfr.bib
% Anexo de cadenas: anexo_cadenas.tex (python Articulos/anexo_cadenas.py)
%
% Los datos que aún no existen se marcan con \pendiente{...} (texto en rojo).
%
% Compilar:  pdflatex rsl ; bibtex rsl ; pdflatex rsl ; pdflatex rsl
%%%%%%%%%%%%%%%%%%%%%%%%%%%%%%%%%%%%%%%%%%%%%%%%%%%%%%%%%%%%%%%%%%%%%%%%%%%%%%%%

\documentclass[conference]{IEEEtran}

\usepackage[utf8]{inputenc}
\usepackage[T1]{fontenc}
\usepackage[spanish]{babel}
\usepackage{cite}
\usepackage{amsmath}
\usepackage{booktabs}
\usepackage{array}
\usepackage{tabularx}
\usepackage{graphicx}
\usepackage{xcolor}
\usepackage{tikz}
\usetikzlibrary{positioning, arrows.meta}
\usepackage[hidelinks]{hyperref}

\addto\captionsspanish{%
  \renewcommand{\tablename}{Tabla}%
  \renewcommand{\figurename}{Fig.}%
}
\renewcommand{\arraystretch}{1.15}
\newcolumntype{Y}{>{\raggedright\arraybackslash}X}
\newcolumntype{L}[1]{>{\raggedright\arraybackslash}p{#1}}
\hyphenation{re-que-ri-mien-tos}

% Datos pendientes (se reemplazan cuando estén disponibles)
\newcommand{\pendiente}[1]{\textcolor{red}{[#1]}}
\newcommand{\kappaFleiss}{\pendiente{por completar}}
\newcommand{\kappaCohen}{\pendiente{por completar}}

\begin{document}

\title{Revisión Sistemática de la Literatura sobre la Clasificación de Requerimientos No Funcionales con Desbalance de Clases: Aumentación de Datos, BERT y Modelos de Lenguaje de Gran Escala}

\author{%
\IEEEauthorblockN{Andres Lino Montes Mamani, Cristian Cabana Sulca, Alessandro Pastor Mamani Mamani}
\IEEEauthorblockA{Escuela Profesional de Ingeniería de Sistemas\\
Universidad Peruana Unión, Juliaca, Perú\\
Asesor: Mg. Esteban Tocto Cano}
}

\maketitle

\begin{abstract}
\textbf{Contexto.} La clasificación automática de requerimientos no funcionales (NFR) está limitada por el desbalance extremo de clases, y la literatura ofrece técnicas de aumentación de datos, clasificadores basados en BERT y modelos de lenguaje de gran escala (LLM) sin compararlos de forma controlada. \textbf{Objetivo.} Identificar, evaluar y sintetizar la evidencia sobre las técnicas contra el desbalance en la clasificación de NFR, los clasificadores con los que se combinan, el uso de LLM como clasificadores, las fuentes de datos y el rigor con que se evalúan. \textbf{Método.} Revisión sistemática según las guías de Kitchenham y Charters y la declaración PRISMA 2020, con búsquedas en Scopus, Web of Science, IEEE Xplore y ACM Digital Library (2017--2026), criterios explícitos de selección, validación del cribado por dos revisores, evaluación de calidad con ocho criterios y extracción de datos del texto completo. \textbf{Resultados.} \pendiente{se completan tras la extracción del texto completo}. \textbf{Conclusiones.} \pendiente{se completan tras la síntesis}.
\end{abstract}

\begin{IEEEkeywords}
revisión sistemática, requerimientos no funcionales, desbalance de clases, aumentación de datos, BERT, modelos de lenguaje de gran escala
\end{IEEEkeywords}

%==============================================================================
\section{Introducción}
%==============================================================================
Los requerimientos no funcionales (NFR) expresan atributos de calidad como la seguridad, el desempeño, la portabilidad o la tolerancia a fallos, y su clasificación automática a partir de las especificaciones de software es una tarea consolidada en la ingeniería de requerimientos~\cite{ClelandHuang2007-NFR,Kurtanovic2017-SupervisedML,Hey2020-NoRBERT}. El principal obstáculo es el desbalance de clases: en la partición de entrenamiento que usa TAOS~\cite{BenAmmar2025-TAOS}, la clase funcional tiene 557 requerimientos y portabilidad, solo dos.

El estado del arte asociado a esta investigación mostró que la aumentación de datos, los clasificadores basados en BERT y los LLM se han estudiado por separado. Cada estudio evalúa una única técnica contra el desbalance, y los LLM se comparan con transformers que no fueron entrenados con datos aumentados. Aquel análisis se basó en los resúmenes de los estudios de una sola base de datos. Esta revisión sistemática amplía la búsqueda a cuatro bases, valida la selección con dos revisores, evalúa la calidad de cada estudio y extrae los datos del texto completo, de modo que las brechas queden respaldadas por evidencia verificable.

La Sección~\ref{sec:protocolo} describe el protocolo, la~\ref{sec:resultados} presenta los resultados por pregunta de investigación, la~\ref{sec:discusion} los discute, la~\ref{sec:amenazas} analiza las amenazas a la validez y la~\ref{sec:conclusiones} concluye.

%==============================================================================
\section{Protocolo de la revisión}\label{sec:protocolo}
%==============================================================================
La revisión sigue las guías para revisiones sistemáticas en ingeniería de software de Kitchenham y Charters~\cite{Kitchenham2007-Guidelines,Kitchenham2015-EBSE} y reporta el proceso conforme a la declaración PRISMA 2020~\cite{Page2021-PRISMA}.

\subsection{Preguntas de investigación}
La Tabla~\ref{tab:picoc} delimita el alcance con el esquema PICOC y la Tabla~\ref{tab:preguntas} formula las preguntas de investigación.

\begin{table}[!t]
\caption{Alcance de la revisión según el esquema PICOC}
\label{tab:picoc}
\centering\footnotesize
\begin{tabularx}{\columnwidth}{@{}L{1.6cm} Y@{}}
\toprule
\textbf{Elemento} & \textbf{Definición} \\
\midrule
Población & Estudios que clasifican automáticamente requerimientos de software escritos en lenguaje natural en tipos de requerimiento o atributos de calidad. \\
Intervención & Técnicas contra el desbalance y la escasez de datos (remuestreo tipo SMOTE y SMOTE-Tomek, aumentación léxica tipo EDA, TextAttack y TAOS, retrotraducción, generación con LLM, funciones de pérdida ponderadas) y clasificadores basados en transformers o en LLM. \\
Comparación & Clasificadores entrenados sin tratamiento del desbalance, aprendizaje automático clásico y transformers ajustados sin aumentación. \\
Resultados & F1 macro, F1 ponderado, F1 por clase y MCC; desempeño en las clases minoritarias; rigor del protocolo de evaluación. \\
Contexto & Ingeniería de requerimientos; conjuntos como PROMISE, PROMISE\_exp, PURE y SecReq, y texto de usuarios o de proyectos de código abierto. \\
\bottomrule
\end{tabularx}
\end{table}

\begin{table}[!t]
\caption{Preguntas de investigación}
\label{tab:preguntas}
\centering\footnotesize
\begin{tabularx}{\columnwidth}{@{}p{0.6cm} Y@{}}
\toprule
\textbf{ID} & \textbf{Pregunta y motivación} \\
\midrule
PR1 & ¿Qué enfoques de clasificación se aplican a los NFR y cómo han evolucionado? \emph{Sitúa el desempeño de BERT y de los LLM frente a los métodos previos.} \\
PR2 & ¿Qué técnicas se usan contra el desbalance y la escasez de datos, con qué clasificadores y con qué efecto en las clases minoritarias? \emph{Núcleo de la comparación entre TAOS, EDA y SMOTE-Tomek.} \\
PR3 & ¿Cómo se usan los LLM como clasificadores de requerimientos y frente a qué se comparan? \emph{Fundamenta la comparación con BERT aumentado.} \\
PR4 & ¿Qué conjuntos y fuentes de datos se emplean y cómo se amplían? \emph{Respalda la ampliación con requerimientos de GitHub.} \\
PR5 & ¿Con qué rigor se evalúan los estudios (partición, fuga de datos, métricas por clase, pruebas estadísticas, reproducibilidad)? \emph{Determina la confianza que merece la evidencia.} \\
\bottomrule
\end{tabularx}
\end{table}

\subsection{Fuentes y cadenas de búsqueda}
Se consultaron cuatro bases de datos: Scopus, Web of Science Core Collection, IEEE Xplore y ACM Digital Library. En cada una se aplicaron cuatro cadenas equivalentes (Tabla~\ref{tab:busquedas}). Todas comparten un bloque común que combina los términos referidos a requerimientos no funcionales, de calidad o de software con el concepto de clasificación, y cada cadena añade uno de los ejes temáticos de la revisión. Donde la base lo permite, se emplean operadores de proximidad para evitar siglas ambiguas como SRS o EDA. Los filtros comunes son el periodo 2017--2026, los idiomas inglés y español, y los tipos artículo, ponencia, revisión y capítulo de libro. Las cadenas completas figuran en el Anexo~A. Las búsquedas en Scopus se realizaron en octubre de 2026 y las de las demás bases, en \pendiente{fecha}.

\begin{table}[!t]
\caption{Registros recuperados por base de datos y eje temático}
\label{tab:busquedas}
\centering\footnotesize
\begin{tabularx}{\columnwidth}{@{}Y rrrr r@{}}
\toprule
\textbf{Base} & \textbf{Aum.} & \textbf{BERT} & \textbf{LLM} & \textbf{Ampl.} & \textbf{Total} \\
\midrule
Scopus & 84 & 214 & 173 & 94 & 565 \\
Web of Science & \pendiente{} & \pendiente{} & \pendiente{} & \pendiente{} & \pendiente{} \\
IEEE Xplore & \pendiente{} & \pendiente{} & \pendiente{} & \pendiente{} & \pendiente{} \\
ACM DL & \pendiente{} & \pendiente{} & \pendiente{} & \pendiente{} & \pendiente{} \\
\midrule
\textbf{Total} & & & & & \pendiente{} \\
\bottomrule
\end{tabularx}
\end{table}

\subsection{Criterios de selección}
La Tabla~\ref{tab:criterios} define los criterios. Un estudio se incluye si cumple los cuatro criterios de inclusión (CI) y se excluye si cumple cualquiera de exclusión (CE). El texto de usuarios (reseñas, solicitudes, issues) solo se incluye cuando las clases son tipos de requerimiento o atributos de calidad. Cuando el título y el resumen no bastan para decidir, el estudio pasa a la lectura del texto completo.

\begin{table}[!t]
\caption{Criterios de inclusión (CI) y exclusión (CE)}
\label{tab:criterios}
\centering\footnotesize
\begin{tabularx}{\columnwidth}{@{}p{0.75cm} Y@{}}
\toprule
\textbf{Código} & \textbf{Criterio} \\
\midrule
CI1 & Estudio primario que clasifica automáticamente texto de requerimientos en tipos de requerimiento (FR/NFR o subclases de NFR) o en atributos de calidad. \\
CI2 & Usa aprendizaje automático, aprendizaje profundo, transformers o LLM. \\
CI3 & Reporta una evaluación empírica cuantitativa. \\
CI4 & Publicado entre 2017 y 2026, en inglés o español. \\
\midrule
CE1 & Fuera del dominio de software: \emph{requirement} o \emph{quality} se usan con otro sentido. \\
CE2 & Otra tarea de ingeniería de requerimientos o de software (ambigüedad, trazabilidad, generación, priorización, riesgo, bugs), o clases que no son tipos de requerimiento. \\
CE3 & Sin evaluación empírica. \\
CE4 & Estudio secundario: se usa como contexto, no en la síntesis. \\
CE5 & Duplicado, o versión resumida o extendida de otro registro. \\
\bottomrule
\end{tabularx}
\end{table}

\subsection{Proceso de selección y fiabilidad}
La selección tiene tres etapas: eliminación de duplicados entre bases y cadenas, cribado por título y resumen, y evaluación de elegibilidad con el texto completo (Fig.~\ref{fig:prisma}). El cribado inicial se realizó con apoyo de un modelo de lenguaje, y cada decisión se documentó junto con el criterio aplicado. Para evaluar su fiabilidad, dos autores revisaron de forma independiente y ciega una muestra formada por todos los estudios incluidos y el 20\,\% de los excluidos. Cada uno trabajó con un orden aleatorio distinto y sin conocer las decisiones del cribado inicial ni las del otro revisor. El acuerdo entre los tres evaluadores, medido con el kappa de Fleiss, fue $\kappa$ = \kappaFleiss, y el kappa de Cohen entre pares osciló en \kappaCohen. Los desacuerdos se resolvieron por mayoría y discusión.

En la fase realizada en Scopus, los 565 registros se redujeron a 434 tras eliminar 131 duplicados entre cadenas. El cribado por título y resumen excluyó 277 (142 por CE1, 129 por CE2, 4 por CE3 y 2 por CE5), identificó 12 estudios secundarios y retuvo 145 estudios primarios para la lectura del texto completo.

\begin{figure}[!t]
\centering
\resizebox{\columnwidth}{!}{%
\begin{tikzpicture}[
  caja/.style={draw, rounded corners=2pt, align=center, font=\footnotesize, text width=3.4cm, inner sep=3pt, fill=blue!6},
  excl/.style={draw, rounded corners=2pt, align=left, font=\scriptsize, text width=3.6cm, inner sep=3pt, fill=gray!8},
  etapa/.style={font=\scriptsize\bfseries, rotate=90, anchor=south},
  arr/.style={-{Latex[length=1.8mm]}, thick}]
  \node[caja] (id) {Registros identificados\\Scopus 565 · WoS \pendiente{}\\IEEE \pendiente{} · ACM \pendiente{}\\$n =$ \pendiente{}};
  \node[caja, below=6mm of id] (cr) {Registros cribados por\\título y resumen\\$n =$ \pendiente{}};
  \node[caja, below=6mm of cr] (tc) {Textos completos\\evaluados\\$n =$ \pendiente{}};
  \node[caja, below=6mm of tc, fill=blue!15] (in) {Estudios incluidos\\en la síntesis\\$n =$ \pendiente{}};
  \node[excl, right=6mm of id] (e1) {Duplicados: \pendiente{}};
  \node[excl, right=6mm of cr] (e2) {Excluidos: \pendiente{}\\(CE1, CE2, CE3, CE5)\\Secundarios (CE4): \pendiente{}};
  \node[excl, right=6mm of tc] (e3) {Sin acceso: \pendiente{}\\Excluidos con texto\\completo: \pendiente{}};
  \draw[arr] (id) -- (cr); \draw[arr] (cr) -- (tc); \draw[arr] (tc) -- (in);
  \draw[arr] (id) -- (e1); \draw[arr] (cr) -- (e2); \draw[arr] (tc) -- (e3);
  \node[etapa, left=2mm of id.west, anchor=south] {Identificación};
  \node[etapa, left=2mm of cr.west, anchor=south] {Cribado};
  \node[etapa, left=2mm of tc.west, anchor=south] {Elegibilidad};
  \node[etapa, left=2mm of in.west, anchor=south] {Inclusión};
\end{tikzpicture}}
\caption{Diagrama de flujo PRISMA 2020 de la selección de estudios.}
\label{fig:prisma}
\end{figure}

\subsection{Evaluación de la calidad}
La calidad de cada estudio incluido se evaluó con ocho criterios (Tabla~\ref{tab:calidad}), adaptados de Dybå y Dingsøyr~\cite{Dyba2008-Agile} y de las guías de Kitchenham y Charters~\cite{Kitchenham2007-Guidelines} al problema de la clasificación con desbalance. Cada criterio se puntúa con 1 (sí), 0.5 (parcialmente) o 0 (no), lo que da un total entre 0 y 8. La puntuación no se usó para excluir estudios, sino para ponderar la evidencia en la síntesis. Dos autores revisaron de forma independiente el 20\,\% de las evaluaciones.

\begin{table}[!t]
\caption{Criterios de evaluación de la calidad}
\label{tab:calidad}
\centering\footnotesize
\begin{tabularx}{\columnwidth}{@{}p{0.65cm} Y@{}}
\toprule
\textbf{ID} & \textbf{Criterio} \\
\midrule
QA1 & Los objetivos o las preguntas de investigación están claramente definidos. \\
QA2 & El conjunto de datos está descrito: origen, tamaño, clases y su distribución. \\
QA3 & El protocolo de evaluación está descrito y evita la fuga de datos: la aumentación se aplica solo al entrenamiento y no hay duplicados entre entrenamiento y prueba. \\
QA4 & Usa métricas adecuadas para el desbalance (F1 macro, F1 por clase o MCC). \\
QA5 & Compara contra una línea base o contra el estado del arte. \\
QA6 & Reporta variabilidad o pruebas de significancia estadística. \\
QA7 & Es reproducible: código, datos e hiperparámetros disponibles o descritos. \\
QA8 & Discute limitaciones o amenazas a la validez. \\
\bottomrule
\end{tabularx}
\end{table}

\subsection{Extracción y síntesis de datos}
Los datos se extrajeron del texto completo con un formulario estandarizado de 23 campos, agrupados en cinco bloques: datos y tarea (fuente del texto, conjunto, idioma, número de requerimientos y de clases, distribución de clases), tratamiento del desbalance (técnica y si la aumentación se aplica solo al entrenamiento), clasificación (clasificadores y papel del LLM), evaluación (protocolo, control de duplicados, métricas, F1 macro, ponderado y de las clases minoritarias, línea base y pruebas estadísticas) y reproducibilidad (disponibilidad de código y datos, y limitaciones declaradas). Dado que los estudios difieren en conjuntos, clases, particiones y métricas, no procede un metaanálisis. La síntesis es narrativa y tabular por pregunta de investigación, y la evidencia se pondera según la puntuación de calidad.

%==============================================================================
\section{Resultados}\label{sec:resultados}
%==============================================================================
\subsection{Estudios seleccionados}
\pendiente{Tras unificar las cuatro bases, describir el número de estudios incluidos, su distribución por año, tipo de publicación y fuente, y la puntuación de calidad.}

\subsection{PR1: enfoques de clasificación}
\pendiente{Síntesis a partir del formulario de extracción.}

\subsection{PR2: técnicas contra el desbalance}
\pendiente{Síntesis a partir del formulario de extracción.}

\subsection{PR3: LLM como clasificadores}
\pendiente{Síntesis a partir del formulario de extracción.}

\subsection{PR4: conjuntos y fuentes de datos}
\pendiente{Síntesis a partir del formulario de extracción.}

\subsection{PR5: rigor de la evaluación}
\pendiente{Síntesis de los criterios QA1--QA8. Se incluirá la verificación de la partición pública de TAOS: el 52\,\% de sus requerimientos de prueba también figura en entrenamiento y sus 1,534 filas corresponden a 988 requerimientos distintos.}

%==============================================================================
\section{Discusión}\label{sec:discusion}
%==============================================================================
\pendiente{Interpretación de los resultados, implicaciones para la investigación propuesta y comparación con el estado del arte.}

%==============================================================================
\section{Amenazas a la validez}\label{sec:amenazas}
%==============================================================================
\textbf{Validez de constructo.} Las cadenas de búsqueda pueden omitir estudios que describen la tarea con otros términos. Para mitigarlo, se combinaron cuatro ejes temáticos, se emplearon sinónimos y variantes ortográficas y se usaron operadores de proximidad.

\textbf{Sesgo de selección.} El cribado inicial se realizó con apoyo de un modelo de lenguaje. Para controlarlo, dos autores validaron de forma independiente y ciega una muestra de las decisiones, se midió el acuerdo con kappa y se documentó el criterio aplicado a cada registro.

\textbf{Sesgo de publicación.} La revisión incluye solo literatura indexada en cuatro bases; pueden faltar estudios con resultados negativos o publicados en otros medios.

\textbf{Validez de la extracción.} La extracción la realizó un revisor con un formulario estandarizado; dos autores revisaron una muestra.

\textbf{Validez de las conclusiones.} La heterogeneidad de conjuntos, clases y métricas impide comparar cifras entre estudios; por ello, la síntesis es narrativa y se pondera por calidad.

%==============================================================================
\section{Conclusiones}\label{sec:conclusiones}
%==============================================================================
\pendiente{Respuesta breve a cada pregunta de investigación y brechas confirmadas.}

\textbf{Disponibilidad de datos.} Las cadenas de búsqueda, las decisiones de cribado con el criterio aplicado, el formulario de extracción y el código del análisis están disponibles en el repositorio del proyecto (\url{https://github.com/AlessandroPastor/research}).

\bibliographystyle{IEEEtran}
\bibliography{referencias_nfr}

\onecolumn
\appendix
\section*{Anexo A. Cadenas de búsqueda completas}

\addcontentsline{toc}{section}{Anexo A. Cadenas de búsqueda completas}

\subsection*{Scopus (Advanced document search)}

\textbf{S1: aumentación y desbalance.}

\begin{flushleft}\ttfamily\scriptsize\detokenize{TITLE-ABS-KEY ( ( ( "non-functional" OR "non functional" OR nonfunctional OR quality OR NFR OR NFRs ) W/3 requirement* ) OR "software requirement*" OR "requirements specification*" OR "requirement* classification" OR "PROMISE NFR" OR "PROMISE_exp" ) AND TITLE-ABS-KEY ( classif* OR categori* ) AND TITLE-ABS-KEY ( "data augmentation" OR "text augmentation" OR "easy data augmentation" OR TextAttack OR TAOS OR "synthetic data" OR "synthetic sample*" OR oversampl* OR "over-sampling" OR resampl* OR SMOTE OR Tomek OR "class imbalance" OR "imbalanced data*" OR "imbalanced class*" OR "minority class*" OR "back translation" OR paraphras* ) AND PUBYEAR > 2016 AND PUBYEAR < 2027 AND ( LIMIT-TO ( DOCTYPE,"ar" ) OR LIMIT-TO ( DOCTYPE,"cp" ) OR LIMIT-TO ( DOCTYPE,"re" ) OR LIMIT-TO ( DOCTYPE,"ch" ) ) AND ( LIMIT-TO ( LANGUAGE,"English" ) OR LIMIT-TO ( LANGUAGE,"Spanish" ) )}\end{flushleft}

\textbf{S2: BERT y transformers.}

\begin{flushleft}\ttfamily\scriptsize\detokenize{TITLE-ABS-KEY ( ( ( "non-functional" OR "non functional" OR nonfunctional OR quality OR NFR OR NFRs ) W/3 requirement* ) OR "software requirement*" OR "requirements specification*" OR "requirement* classification" OR "PROMISE NFR" OR "PROMISE_exp" ) AND TITLE-ABS-KEY ( classif* OR categori* ) AND TITLE-ABS-KEY ( BERT OR RoBERTa OR DistilBERT OR NoRBERT OR transformer* OR "pre-trained language model*" OR "pretrained language model*" OR "transfer learning" OR "fine tun*" ) AND PUBYEAR > 2016 AND PUBYEAR < 2027 AND ( LIMIT-TO ( DOCTYPE,"ar" ) OR LIMIT-TO ( DOCTYPE,"cp" ) OR LIMIT-TO ( DOCTYPE,"re" ) OR LIMIT-TO ( DOCTYPE,"ch" ) ) AND ( LIMIT-TO ( LANGUAGE,"English" ) OR LIMIT-TO ( LANGUAGE,"Spanish" ) )}\end{flushleft}

\textbf{S3: LLM como clasificador.}

\begin{flushleft}\ttfamily\scriptsize\detokenize{TITLE-ABS-KEY ( ( ( "non-functional" OR "non functional" OR nonfunctional OR quality OR NFR OR NFRs ) W/3 requirement* ) OR "software requirement*" OR "requirements specification*" OR "requirement* classification" OR "PROMISE NFR" OR "PROMISE_exp" ) AND TITLE-ABS-KEY ( classif* OR categori* ) AND TITLE-ABS-KEY ( "large language model*" OR LLM OR LLMs OR GPT* OR ChatGPT OR "generative AI" OR Llama OR Gemma OR Mistral OR Claude OR "zero-shot" OR "few-shot" OR "in-context learning" OR prompt* ) AND PUBYEAR > 2016 AND PUBYEAR < 2027 AND ( LIMIT-TO ( DOCTYPE,"ar" ) OR LIMIT-TO ( DOCTYPE,"cp" ) OR LIMIT-TO ( DOCTYPE,"re" ) OR LIMIT-TO ( DOCTYPE,"ch" ) ) AND ( LIMIT-TO ( LANGUAGE,"English" ) OR LIMIT-TO ( LANGUAGE,"Spanish" ) )}\end{flushleft}

\textbf{S4: GitHub y auto-etiquetado.}

\begin{flushleft}\ttfamily\scriptsize\detokenize{TITLE-ABS-KEY ( ( ( "non-functional" OR "non functional" OR nonfunctional OR quality OR NFR OR NFRs ) W/3 requirement* ) OR "software requirement*" OR "requirements specification*" OR "requirement* classification" OR "PROMISE NFR" OR "PROMISE_exp" ) AND TITLE-ABS-KEY ( classif* OR categori* ) AND TITLE-ABS-KEY ( GitHub OR "issue tracker*" OR "issue report*" OR "open source" OR "open-source" OR "semi-supervised" OR "self-training" OR "pseudo-label*" OR "active learning" OR "weak supervision" OR "dataset construction" OR "gold standard" ) AND PUBYEAR > 2016 AND PUBYEAR < 2027 AND ( LIMIT-TO ( DOCTYPE,"ar" ) OR LIMIT-TO ( DOCTYPE,"cp" ) OR LIMIT-TO ( DOCTYPE,"re" ) OR LIMIT-TO ( DOCTYPE,"ch" ) ) AND ( LIMIT-TO ( LANGUAGE,"English" ) OR LIMIT-TO ( LANGUAGE,"Spanish" ) )}\end{flushleft}

\subsection*{Web of Science Core Collection (Advanced search)}

\textbf{W1: aumentación y desbalance.}

\begin{flushleft}\ttfamily\scriptsize\detokenize{TS=((("non-functional" OR "non functional" OR nonfunctional OR quality OR NFR OR NFRs) NEAR/3 requirement*) OR "software requirement*" OR "requirements specification*" OR "requirement* classification" OR "PROMISE NFR") AND TS=(classif* OR categori*) AND TS=("data augmentation" OR "text augmentation" OR "easy data augmentation" OR TextAttack OR TAOS OR "synthetic data" OR "synthetic sample*" OR oversampl* OR "over-sampling" OR resampl* OR SMOTE OR Tomek OR "class imbalance" OR "imbalanced data*" OR "imbalanced class*" OR "minority class*" OR "back translation" OR paraphras*) AND PY=(2017-2026) AND DT=(Article OR "Proceedings Paper" OR Review OR "Book Chapter") AND LA=(English OR Spanish)}\end{flushleft}

\textbf{W2: BERT y transformers.}

\begin{flushleft}\ttfamily\scriptsize\detokenize{TS=((("non-functional" OR "non functional" OR nonfunctional OR quality OR NFR OR NFRs) NEAR/3 requirement*) OR "software requirement*" OR "requirements specification*" OR "requirement* classification" OR "PROMISE NFR") AND TS=(classif* OR categori*) AND TS=(BERT OR RoBERTa OR DistilBERT OR NoRBERT OR transformer* OR "pre-trained language model*" OR "pretrained language model*" OR "transfer learning" OR "fine tun*") AND PY=(2017-2026) AND DT=(Article OR "Proceedings Paper" OR Review OR "Book Chapter") AND LA=(English OR Spanish)}\end{flushleft}

\textbf{W3: LLM como clasificador.}

\begin{flushleft}\ttfamily\scriptsize\detokenize{TS=((("non-functional" OR "non functional" OR nonfunctional OR quality OR NFR OR NFRs) NEAR/3 requirement*) OR "software requirement*" OR "requirements specification*" OR "requirement* classification" OR "PROMISE NFR") AND TS=(classif* OR categori*) AND TS=("large language model*" OR LLM OR LLMs OR GPT* OR ChatGPT OR "generative AI" OR Llama OR Gemma OR Mistral OR Claude OR "zero-shot" OR "few-shot" OR "in-context learning" OR prompt*) AND PY=(2017-2026) AND DT=(Article OR "Proceedings Paper" OR Review OR "Book Chapter") AND LA=(English OR Spanish)}\end{flushleft}

\textbf{W4: GitHub y auto-etiquetado.}

\begin{flushleft}\ttfamily\scriptsize\detokenize{TS=((("non-functional" OR "non functional" OR nonfunctional OR quality OR NFR OR NFRs) NEAR/3 requirement*) OR "software requirement*" OR "requirements specification*" OR "requirement* classification" OR "PROMISE NFR") AND TS=(classif* OR categori*) AND TS=(GitHub OR "issue tracker*" OR "issue report*" OR "open source" OR "semi-supervised" OR "self-training" OR "pseudo-label*" OR "active learning" OR "weak supervision" OR "dataset construction" OR "gold standard") AND PY=(2017-2026) AND DT=(Article OR "Proceedings Paper" OR Review OR "Book Chapter") AND LA=(English OR Spanish)}\end{flushleft}

\subsection*{IEEE Xplore (Command Search)}

Filtros aplicados en la interfaz: años 2017--2026; tipo de contenido: conferencias, revistas, magazines y libros.

\textbf{I1: aumentación y desbalance.}

\begin{flushleft}\ttfamily\scriptsize\detokenize{("non-functional requirement*" OR "nonfunctional requirement*" OR "quality requirement*" OR "software requirement*" OR "requirements specification*" OR "requirements classification") AND (classif* OR categori*) AND ("data augmentation" OR "text augmentation" OR "synthetic data" OR oversampling OR SMOTE OR Tomek OR "class imbalance" OR "imbalanced data" OR TextAttack OR "back translation")}\end{flushleft}

\textbf{I2: BERT y transformers.}

\begin{flushleft}\ttfamily\scriptsize\detokenize{("non-functional requirement*" OR "nonfunctional requirement*" OR "quality requirement*" OR "software requirement*" OR "requirements specification*" OR "requirements classification") AND (classif* OR categori*) AND (BERT OR RoBERTa OR DistilBERT OR transformer* OR "pre-trained language model*" OR "transfer learning" OR "fine-tuning")}\end{flushleft}

\textbf{I3: LLM como clasificador.}

\begin{flushleft}\ttfamily\scriptsize\detokenize{("non-functional requirement*" OR "nonfunctional requirement*" OR "quality requirement*" OR "software requirement*" OR "requirements specification*" OR "requirements classification") AND (classif* OR categori*) AND ("large language model*" OR LLM* OR GPT* OR ChatGPT OR "generative AI" OR "zero-shot" OR "few-shot" OR prompt*)}\end{flushleft}

\textbf{I4: GitHub y auto-etiquetado.}

\begin{flushleft}\ttfamily\scriptsize\detokenize{("non-functional requirement*" OR "nonfunctional requirement*" OR "quality requirement*" OR "software requirement*" OR "requirements specification*" OR "requirements classification") AND (classif* OR categori*) AND (GitHub OR "issue report*" OR "open source" OR "semi-supervised" OR "self-training" OR "pseudo-label*" OR "active learning")}\end{flushleft}

\subsection*{ACM Digital Library (Advanced Search, campo \emph{Abstract})}

Las tres filas se combinan con AND. Filtro de fecha de publicación: 2017--2026.

\textbf{Fila 1 (común):}

\begin{flushleft}\ttfamily\scriptsize\detokenize{"non-functional requirements" OR "non-functional requirement" OR "nonfunctional requirements" OR "quality requirements" OR "software requirements" OR "requirements classification" OR "requirements specification"}\end{flushleft}

\textbf{Fila 2 (común):}

\begin{flushleft}\ttfamily\scriptsize\detokenize{classification OR classify OR classifying OR categorization OR categorizing}\end{flushleft}

\textbf{A1 --- Fila 3: aumentación y desbalance.}

\begin{flushleft}\ttfamily\scriptsize\detokenize{"data augmentation" OR "text augmentation" OR "synthetic data" OR oversampling OR SMOTE OR Tomek OR "class imbalance" OR "imbalanced data" OR TextAttack OR "back translation"}\end{flushleft}

\textbf{A2 --- Fila 3: BERT y transformers.}

\begin{flushleft}\ttfamily\scriptsize\detokenize{BERT OR RoBERTa OR DistilBERT OR transformer OR transformers OR "pre-trained language model" OR "pretrained language model" OR "transfer learning" OR "fine-tuning"}\end{flushleft}

\textbf{A3 --- Fila 3: LLM como clasificador.}

\begin{flushleft}\ttfamily\scriptsize\detokenize{"large language model" OR "large language models" OR LLM OR LLMs OR GPT OR ChatGPT OR "generative AI" OR "zero-shot" OR "few-shot" OR prompting}\end{flushleft}

\textbf{A4 --- Fila 3: GitHub y auto-etiquetado.}

\begin{flushleft}\ttfamily\scriptsize\detokenize{GitHub OR "issue reports" OR "open source" OR "semi-supervised" OR "self-training" OR "pseudo-labeling" OR "active learning"}\end{flushleft}


\end{document}
